\section{Changes}
\label{app:changes}

Decisions taken about the paper's framing, recorded here so they are not relitigated and so a
co-author can see what was deliberate. Not part of the paper.

\subsection{The free-verification framing is removed}

\textbf{Decision.} The phrase \emph{free exact verification}, and the framing built on it, is out.
It is gone from the title, from the abstract, and from the body. Where the property still matters
it is stated as what it is: the judge executes the model's proposal, and the engine that generates
a sample is the engine that rules on a step.

\textbf{Why.} Free verification is not a distinguishing property. It is the ordinary condition of
the benchmarks this paper sits beside: code benchmarks run tests, maths benchmarks run a proof
assistant, and nobody advertises that as a contribution because it is assumed. Leading with it
invites two bad readings at once. A reviewer who works on code or maths benchmarks reads it as a
claim to something everyone already has, which reads as naivety. A reviewer who does not reads it
as the paper's main idea, which displaces the actual contribution. Either way the reader's
attention goes to the verifier rather than to the dataset, the task and the tiers, which is where
the work is.

\textbf{What replaces it.} The contribution is stated as a dataset and an evaluation framework for
spatial and geometric reasoning. Exact verification stays in the paper as a property of the
environment (Section 5) and as the reason the scoring protocol can be trusted (Section 6.4), which
is the right altitude for it: a mechanism the paper relies on, not the thing it is selling.

\subsection{The task definition is fixed and is now stated once, precisely}

\textbf{Decision.} The model is given \textbf{both} the crease pattern \textbf{and} the final
folded state, and is asked for the sequence of folds that turns the flat sheet into that state.
Both endpoints are supplied; the path between them is not.

\textbf{Why.} The abstract previously said the task was recovering the sequence ``that produced a
crease pattern,'' which omits the final state and describes a harder and different task. Section 1
and the abstract now agree and use the same wording. Any future edit that changes what the model is
given has to change both.

\subsection{Geometry versus search is not a contradiction, and the resolution is now written down}

\textbf{Decision.} The paper keeps saying that what it scores is search, and also keeps positioning
itself as a benchmark for spatial and geometric reasoning. Both are true and the bridge between
them is now stated in Section 1 and the abstract rather than left implicit.

\textbf{Why.} The two only look like they collide. Geometry is supplied to the model as input, so
the paper is not scoring whether a model can compute a reflection. What it scores is whether a
model can \emph{act} on the geometry it was given in order to choose a direction. That distinction
is load-bearing, because blind search does not survive this task: the reachable space grows with
depth faster than exhaustive strategies can cover, so a model proposing without reading the pattern
exhausts its budget on the deeper strata, and a deterministic breadth-first search fails on the
same samples for the same reason. Blind enumeration is not a weak baseline here, it stops being
viable. A model that cannot read a reflection off a pattern therefore has no way to choose a
direction, and on this task that is indistinguishable from being unable to do it at all. Geometric
understanding is the precondition for the search that gets scored.

\subsection{Ablations are demoted to tiers, because the full study has not been run}

\textbf{Decision.} The word \emph{ablation} is out of the abstract and Section 1. Section 9 is now
``Progressive tiers of tool assistance.'' The four arms are unchanged; only the claim about them is
weaker.

\textbf{Why.} A full ablation study was planned and has not been run. An ablation is a specific
claim, that one factor was isolated and its contribution measured, and a partial sweep does not
support it. Reporting tiers says what was actually done: the model is run at four levels of
assistance and the levels are compared. If the full study is run before submission this can be
restored to an ablation, and Section 9 and the abstract both change back together.

\subsection{Still open}

\begin{itemize}
\item \textbf{The abstract's lede.} It now opens on multimodal models and the geometry gap,
  matching Section 1. It is still long, at roughly one page-third, and should be cut once
  Section 10 exists and the headline result has to fit.
\item \textbf{The headline result sentence} in Section 1 and the abstract. \texttt{[TO RUN]}
\item \textbf{Figure 1}, the generation-versus-recovery asymmetry. Section 12 has it as ``to
  draw''; an introduction to a benchmark paper is usually carried by that figure.
\item \textbf{The star example in Section 1 is an unverified citation} and blocks submission. See
  Section 15 rule 7.
\end{itemize}

\subsection{The 89.3\% scope claim is now bounded to the corpus it was measured on}

\textbf{Decision.} Every statement of the figure now reads ``89.3\% of the 366 real crease patterns
in Flat-Folder's \texttt{examples/instagram/} corpus'' rather than ``89.3\% of real crease
patterns.'' Three places: Section 1 contribution 5, Section 7, Section 11.

\textbf{Why.} The measurement was made on one convenience corpus of 366 patterns. ``Real crease
patterns'' names a population that was never sampled, and a reviewer who checks the provenance will
read the generalization as either careless or deliberate. The bounded version is unattackable and
loses none of its force, because the point is only that the action space excludes most real origami
and the boundary is reported rather than buried. This was the single easiest reviewer objection in
the paper to remove.

\subsection{The BFS claim was removed until the number exists}

\textbf{Decision.} Section 1 previously asserted that a deterministic breadth-first search fails on
the same samples a model does. That sentence is gone. In its place is a statement that the depth at
which exhaustive search stops being affordable is measurable, and a \texttt{[TO RUN]} for the
baseline.

\textbf{Why.} The claim had no measurement behind it, and the repo's own baseline suggests it may
point the other way. \texttt{workspace/search\_baseline.mjs} was written to answer the opposite
worry: its header records that \texttt{list\_legal\_folds} evaluates roughly 2010 candidate actions
per state and returns a mean of 3.51 that are legal and stay inside the target CP, measured over
73,750 enumerations, and warns that if exhaustive search over that pruned set solves the easy tier
outright then the enumerator arm's 20 to 26 percent is not measuring origami reasoning. A branching
factor of 3.5 is small. Asserting that blind search collapses, while shipping a script that
suspects it does not, is the kind of contradiction a reviewer finds by reading the repository.
\textbf{The script has never had its results written into any document in this project.} Run it and
the sentence can come back, stronger, with a number.

\subsection{Baselines are named, because a benchmark without them is usually rejected}

\textbf{Decision.} Section 8.5 names three: uniform random over enumerated legal actions,
deterministic breadth-first over the same action set, and the symbolic CP$\rightarrow$Seq solver
(Creasy). All \texttt{[TO RUN]}.

\textbf{Why.} Reviewers of dataset and benchmark papers ask what the floor is and what a non-learned
method achieves. Without that, model numbers are uninterpretable: 25 percent is either impressive
or embarrassing depending on what search alone does. The BFS baseline doubles as the answer to the
``the simulator did all the work'' objection of Section 7, which is currently argued structurally
and would be much stronger argued numerically.

\subsection{Licence and release terms}

\textbf{Decision.} Everything ships under \textbf{BSD 3-Clause}: generator, environment and
verifier, corpus checks, comparator, harness, and the corpus itself. Recorded in Section 13a.

\textbf{Why.} Dataset and benchmark tracks expect an explicit licence, a hosting story and a
maintenance commitment, and their absence is a routine reviewer complaint. The corpus shipping as a
seed rather than a blob answers hosting almost for free. Still outstanding: the anonymised
repository URL, an archival DOI, and a datasheet.

\subsection{OrigamiBench is closer to this work than Section 2 assumes, and this needs a decision}

\textbf{Status: open. This is the most significant finding of the September 2026 citation check.}

The missing \texttt{origamibench} citation was recovered: Agarwal, Wu, Jian, Hu, Mansoor, Li, Peng,
Dai, Ding and Sansone, \emph{OrigamiBench: An Interactive Environment to Synthesize Flat-Foldable
Origamis}, arXiv:2603.13856, March 2026. Its abstract describes ``an interactive benchmark in which
models iteratively propose folds and receive feedback on physical validity and similarity to a
target configuration.''

That is much nearer this paper's design than the draft assumed when it listed OrigamiBench as one
of the works that ``complete the picture.'' The claim in Section 1 and Section 2.4 that
\emph{none} of the existing benchmarks steps an engine and rules on a proposed move is now doubtful
as stated, because OrigamiBench appears to do exactly that.

Two distinctions probably survive and both need checking against the paper rather than the
abstract. First, OrigamiBench reports ``similarity to a target configuration'', which is a graded
comparison, where this work replays and compares exactly under a declared symmetry group. Second,
OrigamiBench synthesises toward a target, where this work recovers a sequence from a crease pattern
and a final state, which is the inverse problem. If both hold, the contribution stands and the
wording needs narrowing from ``none of them executes'' to something precise about exactness and
about the inverse direction. \textbf{If neither holds, the framing of Section 1 and Section 2.4 has
to change.} Read the paper before submission.

\subsection{Citations verified in the September 2026 check}

\begin{itemize}
\item \textbf{OrigamiSpace.} 350 instances and the four tasks confirmed. The venue is confirmed as
  \textbf{NeurIPS 2025}, not an unvenued preprint, so Section 15 rule 4 is now satisfied and can be
  retired.
\item \textbf{GamiBench.} 186 regular and 186 impossible patterns, six viewpoints, three VQA tasks,
  and the viewpoint-consistency and impossible-fold-selection-rate metrics all confirmed. Authors
  confirmed.
\item \textbf{OrigamiBench.} Recovered; see Section 16.10.
\item \textbf{COrigami.} Recovered as \emph{COrigami: An AI Pipeline for Co-Designing Flat-Foldable
  Visually Recognisable Origami}, arXiv:2606.26299. Author list still unverified.
\item \textbf{Still unverified:} the \emph{Vision Language Models Are Blind} star example of
  Section 1, and the author lists of thirteen theory and tools entries, which now carry sort keys so
  the bibliography renders but still need real authors.
\end{itemize}

\subsection{Two claims about exactness were narrowed to what the code does}

\textbf{Decision.} The paper no longer calls the pipeline tolerance-free, and no longer describes
the Level 2 group as the eight symmetries of the square plus a translation.

\textbf{Why.} Neither matched\\
\texttt{DHEERAJ\_WORKSPACE/EXPERIMENT\_SETUP/terminal\_match.mjs}.
Two mismatches:

\begin{enumerate}
\item \textbf{The group is larger than claimed.} The comparator quotients out the full plane
  isometry group, including \emph{arbitrary} rotation, not the eight symmetries of the square. Its
  own header says so: ``terminalMatch quotients the plane isometry group out.''
\item \textbf{There is a tolerance.} \texttt{MATCH\_TOL = 2e-6}, with degenerate-vertex cleanup at
  \texttt{1e-7} and a collinearity threshold at \texttt{1e-9}. The value is recorded in every result
  the harness writes: 1,187 of 1,211 attempts across the saved runs carry
  \texttt{"tolerance": 2e-06}.
\end{enumerate}

Admitting arbitrary rotation makes a tolerance unavoidable, because a rotation by an angle that is
not a multiple of a right angle produces irrational coordinates. So the tolerance is not
sloppiness; it is the necessary consequence of a design choice the paper wanted. What was wrong was
the claim, not the code.

The formulation now used: the \textbf{verifier} is exact, because legality is decided
combinatorially by the engine and involves no threshold; \textbf{terminal-state equality under
transforms} carries a stated numerical tolerance. Level 1 crease-set comparison remains genuinely
tolerance-free.

Claiming an exactness the implementation does not have is the first thing a reviewer with the
repository checks, and in a paper whose contribution is the trustworthiness of its judge it would
be the most damaging possible thing to get wrong.

\subsection{Figures have written briefs}

\textbf{Decision.} \texttt{paper/figures/figure1.md} through \texttt{figure6.md}, one per figure,
each naming where it appears, the argument it carries, what it must show, \textbf{what it must not
show}, how to tell whether the drawing worked, and a draft caption.
\texttt{paper/figures/README.md} indexes them.

\textbf{Why.} Three of the six figures can be drawn wrongly in ways that would contradict the
paper's own claims, and the briefs say so explicitly. Figure 1 must not imply legal moves are rare,
which Section 3.3 exists to forbid. Figure 4 must draw the environment and the verifier as one box,
because that is the claim. Figure 6 must not hide a baseline crossing. A figure brief that only says
what to draw would let all three happen.

\subsection{The abstract is rewritten as one paragraph and made consistent with the code}

\textbf{Decision.} The abstract is now a single paragraph of roughly 440 words, down from six
paragraphs, and every claim in it matches the implementation and the rest of the paper.

\textbf{Why one paragraph.} The ICLR template states the requirement outright: ``The abstract must
be limited to one paragraph.'' Six paragraphs was a hard format violation, not a style preference.

\textbf{What changed beyond the format.} Five inconsistencies were removed rather than reworded.

\begin{enumerate}
\item \textbf{The task.} It now says the model is given the crease pattern \emph{and the final
  folded state}, matching Section 1. The old text said only ``the sequence of folds that produced a
  crease pattern'', which describes a harder and different task.
\item \textbf{The dataset.} It now leads with a benchmark \emph{and dataset} of 600 samples. The old
  abstract named no size and made no dataset claim, which for a datasets-and-benchmarks submission
  is the first thing a reviewer looks for.
\item \textbf{Tolerance.} It no longer claims tolerance-free replay. Crease-set comparison is
  tolerance-free and says so; folded-state comparison is up to the plane isometry group and carries
  a stated numerical tolerance, per Section 16.12.
\item \textbf{Ablation vocabulary.} ``Ablations that remove visual feedback, remove filtering,
  remove tools'' is replaced by progressive tiers of tool assistance, per Section 16.4.
\item \textbf{Baselines.} The abstract now says models are ranked against a deterministic search
  baseline rather than only against each other, which is what Section 8.4 promises and what a
  benchmark paper is expected to provide.
\end{enumerate}

The free-verification framing stays out (Section 16.1). What survives of it is the one sentence that
carries a fact rather than a slogan: the verdict on a step carries no threshold because legality is
decided combinatorially.

\textbf{Done, see Section 16.18.} The abstract was cut to 282 words once Section 10 produced
numbers.

\subsection{Figure drafting is scripted, with the limits of the method written into the script}

\textbf{Decision.} \texttt{paper/figures/gen\_figures.sh} drives the imagegen CLI to produce tracing
drafts. \texttt{--list} prints how each figure should actually be made; \texttt{--dry-run} needs no
API key.

\textbf{How it runs.} Each figure is one \texttt{codex exec} call. Codex uses its \textbf{built-in
\texttt{image\_gen} tool}, which is the imagegen skill's own preferred path and needs \textbf{no
\texttt{OPENAI\_API\_KEY}} and no direct API access. The \texttt{scripts/image\_gen.py} CLI
fallback, which does require a key, is explicitly not used.

\textbf{Why only two of six by default.} Figures 1 and 4 are conceptual and draft usefully from a
prompt; a bare run drafts only those. \texttt{--all} overrides it. The other four carry data or
geometry and must not ship as generated images: figures 2 and 5 should be built from renders that
already exist in this repo (\texttt{initial/cp.png}, \texttt{final/top.png},
\texttt{final/exploded.png} for any sample), and figures 3 and 6 must be plotted from the release
manifest and the aggregation output. A generated crease pattern would be a fabricated figure in a
paper whose subject is exact verification, which is the worst place in the literature to put one.
The imagegen skill says the same thing in its own terms: diagrams are ``better produced directly in
SVG, HTML/CSS, or canvas''. The script prints this rather than assuming it is remembered.

\subsection{Draft figures are labelled as drafts, in the PDF itself}

\textbf{Decision.} Every placeholder figure renders a visible banner reading \textbf{DRAFT IMAGE:
PLACEHOLDER, NOT FINAL}, with a line saying it will be replaced by a human-authored figure and that
drafts come from OpenAI \texttt{gpt-image-2}. Captions carry a \texttt{[DRAFT IMAGE]} prefix so the
marking also appears in any list of figures. The same banner is at the top of each
\texttt{paper/figures/figureN.md} and the README.

\textbf{Why in the PDF and not only in the notes.} A placeholder that is only labelled in a side
file becomes a real figure the moment somebody exports a PDF to show a collaborator. The label has
to travel with the artifact. It costs nothing and removes a whole class of accident.

\subsection{LLM usage is disclosed, per ICLR 2026 policy}

\textbf{Decision.} A new Section 13b discloses every use: models as the object of study, Claude Code
for the implementation, partial AI assistance for literature search alongside Google Scholar and
alphaXiv, AI assistance in drafting, and \texttt{gpt-image-2} for draft figures.

\textbf{The policy, checked rather than assumed.} ICLR 2026 has two policies on LLM use. Policy 1:
any use of an LLM must be disclosed, on the Code of Ethics principle that all contributions to the
research must be acknowledged. Policy 2: authors are ultimately responsible for their contributions
and must not make false or misleading claims. The policy does \textbf{not} distinguish a model used
as a research instrument from one used for writing; both must be disclosed.

\textbf{Caveat. Disclosure is required in the paper's text \emph{and} in the submission form.}
Section 13b satisfies the first. The second is a separate field on the OpenReview form and is easy
to miss.

The FAQ does not specify a required section, required wording, or whether the disclosure counts
against the page limit. It is written as a full section here on the basis that under Policy 2 the
authors carry responsibility either way, so under-disclosing buys nothing and risks everything.

\textbf{Why it is written long rather than minimal.} A disclosure that omits a use is worse than one
that reports a use a reader would have forgiven. The section also names the two verification habits
that exist \emph{because} of the reliance on AI-assisted implementation, the independent corpus
checks and the comparator's negative control, and states that the negative control caught a real
defect. That is the honest form of the disclosure: not a claim that assistance introduced no risk,
but a description of what was put in place to catch it.

\subsection{The abstract is cut to 282 words, and the result now carries its ending}

\textbf{Decision.} The abstract is 282 words, down from 443. It ends on the finding rather than on
the scoring protocol.

\textbf{Why it was too long.} 443 words is roughly a page-third and reads as a summary of the
paper's mechanism rather than of its contribution. The cause was structural: the abstract was
written before Section 10 existed, so the middle of it carried mechanism description that was doing
the work a result would otherwise do. Once the numbers existed, most of that description could go.

\textbf{What was cut, and on what principle.} Everything that a reader can reach in one page of
Section 1 and that is not a claim. Gone: the enumeration of refusal types beyond one example, the
two named levels of equality, the crease-set versus folded-state distinction, the per-stratum
reporting protocol, queries-to-solution as a secondary measure, the stratification axes, and the
tolerance discussion. None of that is wrong and all of it survives in Sections 5, 6 and 8; none of
it belongs in an abstract competing with a result for the same 250 words.

\textbf{What was kept, and why each earns its place.} The multimodal framing, because it is the
reason a reader at ICLR should care. The task in one sentence. The generate-versus-solve asymmetry
with the NP-hardness, because it is the paper's foundation. Constructed ground truth and
seed-rebuilding, because those are the dataset claims a benchmark reviewer reads for.
Environment-as-verifier in two sentences. The tiers, in one clause. And the three numbers: no solve
past easy, 54.7 against 20.4, 48 percent cycling.

\textbf{The ending changed, and this is the part that matters most.} The old abstract ended on the
scoring protocol, which tells a reader what we did. It now ends on state cycling, which tells them
what we found. An abstract whose last sentence is a method is a proposal; one whose last sentence is
a finding is a result.

\subsection{No repository or dataset links in the paper}

\textbf{Decision.} Section 13a states the licence and one sentence: code and data are released on
publication, links omitted. No URLs, no DOI, no datasheet, no hosting or maintenance prose.

\textbf{Why no links at all.} Two reasons, and either is sufficient on its own. A URL naming the
authors' organisation identifies them in a double-blind submission. And a link committed to in a
paper is fixed at the moment of submission, so naming one before the repository and dataset are
actually published creates a name the project then has to honour. Omitting them removes both
problems and costs nothing at submission time.

\textbf{When to add them.} At camera-ready, once the repository and the dataset exist under names
that are not going to change. At that point the paper is the authority and the artifacts are renamed
to match it, rather than the paper being edited to chase them.
